\section{Threat Model and Scope}
\label{sec:threat}

We consider a smart home deployment in which IoT devices \textcolor{blue}{obtain software or firmware updates from vendor servers, either directly or through a companion app or desktop application.
An update involves up to three exchanges: an \emph{update check} that returns metadata (the available version and where to get it), an optional \emph{authorization} step in which a vendor service approves installation on this particular device, and the \emph{image transfer}.
Each exchange may use a different server and protocol, and the image may reach the device over the Internet or across the home network from an application that downloaded it.}

\paragraph{\textcolor{blue}{Protection layers}}
\textcolor{blue}{Update security rests on three layers:
(i)~the \emph{channel}, i.e., encryption and server authentication of each exchange;
(ii)~the \emph{image}, i.e., the integrity mechanism inside the update file or its manifest, ranging from checksums such as CRC32 or MD5, which anyone can recompute, to vendor signatures; and
(iii)~\emph{device-side checking}, i.e., whether the device verifies (ii) before it installs the image.
We measure the channel for all devices and characterize the image protection for the images we can obtain; device-side checking cannot be observed from the network and is the subject of firmware analyses such as~\cite{294541}.}

\paragraph{Adversary profile}
We assume a network-level adversary ($\mathcal{A}$) who can observe traffic on any segment of the update path.
Specifically, $\mathcal{A}$ may be (i)~an on-path Wi-Fi eavesdropper within range of the home access point, (ii)~a compromised or backdoored home gateway or router, or (iii)~a transit-level observer (e.g., a compromised ISP node or a state-level actor).
A \emph{passive} $\mathcal{A}$ reads packet metadata (servers, sizes, timing) and any unencrypted content; an \emph{active} $\mathcal{A}$ can also drop, modify, and inject packets.

\paragraph{Adversary goals and what they require}
\emph{Reconnaissance}: learn a device's model, firmware version, and update timing, to find unpatched targets; unencrypted update checks or downloads suffice.
\emph{Tampering}: substitute or modify an image or its metadata in transit; this succeeds only if the channel does not authenticate the server \emph{and} the image is unsigned, or the device does not check its signature.
\emph{Server impersonation}: present a certificate that the device accepts, e.g., by factoring a weak key or by compromising a long-lived vendor certificate authority; it helps $\mathcal{A}$ only for the same unprotected images.
\emph{Rollback and freeze}: replay older images or withhold updates; signed and fresh metadata prevents this, TLS alone does not~\cite{samuel2010tuf}.
Keeping the image secret is not a goal: vendors typically publish their images, so $\mathcal{A}$ can simply download them.

\paragraph{\textcolor{blue}{What the network can and cannot show}}
\textcolor{blue}{In TLS~1.2 the Finished messages authenticate the whole handshake, so an on-path adversary cannot remove strong cipher suites from the list that a client offers; weak \emph{offers} become exploitable only if a server selects them, through export-grade and anonymous suites~\cite{beurdouche2015freak,adrian2015logjam}, or through insecure version fallback~\cite{rfc7507}.
We therefore treat the \emph{negotiated} settings as the actual exposure and the offered suites as legacy support.
TLS~1.3 encrypts the server certificate~\cite{rfc8446}, so for TLS~1.3 connections we see the server name and the negotiated settings but not the certificate.
We can see neither manifests carried inside TLS nor whether a device checks signatures or validates certificates.}

\paragraph{Trust assumptions}
We assume that the vendor's update server itself is not compromised.
We further assume that the IoT device firmware, at the time of deployment, is authentic. The security of the update \emph{delivery channel} between the server and the device is \emph{not} assumed and is the subject of our analysis.
